\documentclass[11pt]{article}
\usepackage[margin=1in]{geometry}
\usepackage[T1]{fontenc}
\usepackage{amsmath,amssymb,mathtools}
\usepackage{booktabs,array,enumitem,caption,float}
\usepackage{xcolor,hyperref}
\sloppy
\hypersetup{colorlinks=true, linkcolor=blue!50!black, urlcolor=blue!50!black}
\newcommand{\code}[1]{\texttt{#1}}
\newcommand{\file}[1]{\texttt{#1}}
\setlist{nosep, leftmargin=1.4em}

\title{AutoBayes against brms: where the speed comes from}
\author{Benchmark run of 2026-10-25, prepared for Andrew}
\date{2026-10-25}

\begin{document}
\maketitle

\noindent\textbf{The answer.} Fitting each model from scratch (compiling and sampling), brms took 0.8 to 18.1 times as long as AutoBayes (median 6.5). Its sampler alone (warm-up and sampling of every chain, compilation excluded) took 0.6 to 13.2 times as long (median 3.8). Splitting that ratio into its four factors (Section~\ref{sec:decomp}), the largest was the gradient factor (for AutoBayes) on 15 models, the overhead factor (against AutoBayes) on 2 models, the gradient factor (against AutoBayes) on 1 model. brms's sampler was the faster on B8, B12.

\section{What was run}

Each of the 18 models was fitted by both libraries with the same formula, family and explicit priors, on the same simulated data (\file{examples/brms\_parity/benchmark\_data.py}, fixed seeds), with the same sampler settings: 4 chains in parallel, 1000 warm-up and 1000 kept draws each, target acceptance 0.9.
\begin{itemize}
\item \textbf{brms} 2.23.0 with rstan 2.32.7 (R 4.5.2), 4 cores; compiled with \code{brm(chains = 0)} and sampled with \code{update(recompile = FALSE)}.
\item \textbf{AutoBayes} (commit \code{ce04dd3}) with NumPyro 0.21.0, JAX 0.10.2 and PyMC 5.28.5 on the CPU (macOS-26.6.1-arm64-arm-64bit, 10 cores), \code{ab.fit} unchanged and instrumented from outside (\file{examples/brms\_parity/README.md}).
\end{itemize}

\section{Time to a posterior, by stage}

\begin{table}[H]
\footnotesize
\setlength{\tabcolsep}{4pt}
\begin{center}
\begin{tabular}{@{}llrrrrrr@{}}
\toprule
& model & rows & codegen & compile & exec & around & total \\
\midrule
B1 & Bernoulli, \code{(1 + x | g)} & 2,000 & 0.09 & 31.1 & 4.5 & 1.2 & 36.8 \\
B2 & Gaussian, \code{(1 + x | g)} & 20,000 & 0.12 & 29.2 & 74.9 & 1.8 & 105.9 \\
B3 & cumulative, five categories & 1,500 & 0.07 & 26.1 & 32.4 & 0.44 & 58.9 \\
B4 & zero-inflated Poisson & 2,000 & 0.09 & 25.5 & 14.6 & 0.49 & 40.6 \\
B5 & Weibull & 2,000 & 0.08 & 26.3 & 5.9 & 0.51 & 32.7 \\
B6 & zero-inflated neg.\ binomial & 2,000 & 0.08 & 26.2 & 15.9 & 0.53 & 42.7 \\
B7 & sratio, five categories & 1,500 & 0.07 & 25.9 & 17.1 & 0.47 & 43.5 \\
B8 & Gaussian, horseshoe on 20 & 400 & 0.09 & 25.5 & 1.7 & 0.73 & 27.9 \\
B9 & beta-binomial, trials & 2,000 & 0.08 & 24.7 & 23.7 & 0.65 & 49.1 \\
B10 & Gaussian, \code{(1 | mm(g1, g2))} & 1,000 & 0.08 & 24.7 & 2.3 & 0.58 & 27.6 \\
B11 & mixture of two Gaussians & 1,000 & 0.10 & 27.0 & 4.0 & 0.51 & 31.5 \\
B12 & exact GP, \code{gp(x)} & 200 & 0.07 & 27.6 & 468.8 & 1.2 & 497.6 \\
B13 & AR(1), \code{ar(t, g)} & 2,000 & 0.10 & 25.0 & 2.1 & 1.2 & 28.3 \\
B14 & non-linear, \code{a * exp(b * x)} & 1,000 & 0.11 & 25.2 & 5.7 & 0.66 & 31.6 \\
B15 & \code{mvbind(y1, y2)}, rescor & 1,000 & 0.17 & 29.7 & 13.1 & 0.88 & 43.8 \\
B16 & monotonic \code{mo(edu)} & 1,500 & 0.12 & 28.2 & 6.5 & 1.1 & 35.7 \\
B17 & Weibull, \code{cens()} & 2,000 & 0.10 & 28.7 & 5.6 & 1.3 & 35.6 \\
B18 & distributional, \code{sigma \textasciitilde{} z} & 2,000 & 0.11 & 26.9 & 8.5 & 0.72 & 36.2 \\
\bottomrule
\end{tabular}\end{center}
\caption{\small brms's stages, seconds. \emph{compile}: \code{brm(chains = 0)}, code generation, C++ compilation and loading; \emph{codegen}: the code generation alone, run again afterwards; \emph{exec}: the slowest chain's warm-up and sampling as Stan reports them; \emph{around}: the rest of \code{update()} (forking the chains, brms's data preparation and its renaming of the fit).}\label{tab:brms}
\end{table}
\begin{table}[H]
\footnotesize
\setlength{\tabcolsep}{4pt}
\begin{center}
\begin{tabular}{@{}lrrrrrrrrr@{}}
\toprule
& build & pre & compile & JAX & exec & post & total & \multicolumn{2}{c}{faster by} \\
\cmidrule(lr){9-10}& & & & & & & & total & exec \\
\midrule
B1 & 2.1 & 0.33 & 1.2 & 1.2 & 2.5 & 0.83 & 7.0 & 5.27 & 1.78 \\
B2 & 10.4 & 0.61 & 1.3 & 1.3 & 17.4 & 0.86 & 30.5 & 3.47 & 4.31 \\
B3 & 0.10 & 0.38 & 1.0 & 1.0 & 4.9 & 0.29 & 6.7 & 8.82 & 6.64 \\
B4 & 0.12 & 0.48 & 1.1 & 1.1 & 4.1 & 0.22 & 6.0 & 6.82 & 3.61 \\
B5 & 0.12 & 0.28 & 1.1 & 1.1 & 4.3 & 0.14 & 6.0 & 5.48 & 1.37 \\
B6 & 0.32 & 0.53 & 1.4 & 1.4 & 12.7 & 0.15 & 15.2 & 2.81 & 1.25 \\
B7 & 0.09 & 0.32 & 1.0 & 1.0 & 5.6 & 0.13 & 7.1 & 6.12 & 3.08 \\
B8 & 0.09 & 0.48 & 1.2 & 1.2 & 3.0 & 0.32 & 5.1 & 5.46 & 0.56 \\
B9 & 0.13 & 0.29 & 1.6 & 1.6 & 20.4 & 0.15 & 22.6 & 2.17 & 1.16 \\
B10 & 0.09 & 0.24 & 1.0 & 1.00 & 0.47 & 0.12 & 1.9 & 14.38 & 4.86 \\
B11 & 0.04 & 0.75 & 1.1 & 1.1 & 0.98 & 0.19 & 3.1 & 10.12 & 4.05 \\
B12 & 0.05 & 0.28 & 1.1 & 1.1 & 617.9 & 0.50 & 619.8 & 0.80 & 0.76 \\
B13 & 0.09 & 0.23 & 0.96 & 0.95 & 0.22 & 0.07 & 1.6 & 18.11 & 9.79 \\
B14 & 0.05 & 0.23 & 0.99 & 0.99 & 0.83 & 0.11 & 2.2 & 14.35 & 6.89 \\
B15 & 0.16 & 0.62 & 1.6 & 1.6 & 1.00 & 0.23 & 3.6 & 12.17 & 13.16 \\
B16 & 0.49 & 0.43 & 1.2 & 1.2 & 0.98 & 0.26 & 3.4 & 10.64 & 6.59 \\
B17 & 0.14 & 0.48 & 1.2 & 1.1 & 4.9 & 0.15 & 6.8 & 5.21 & 1.13 \\
B18 & 0.63 & 0.25 & 1.1 & 1.0 & 1.6 & 0.15 & 3.7 & 9.90 & 5.43 \\
\bottomrule
\end{tabular}\end{center}
\caption{\small AutoBayes's stages, seconds. \emph{build}: the model; \emph{pre}: the fit up to NumPyro (PyTensor's graph turned into JAX, initial points); \emph{compile}: the NumPyro run up to the end of its largest XLA compilation, with \emph{JAX} the part of it JAX's monitoring reports as tracing, lowering and compiling; \emph{exec}: the rest of the run, every chain's warm-up and sampling; \emph{post}: the draws moved to the constrained scale and the InferenceData. The last two columns are brms's time over AutoBayes's, end to end and for the samplers' execution.}\label{tab:autobayes}
\end{table}

\section{The samplers' work}

\begin{table}[H]
\footnotesize
\setlength{\tabcolsep}{4pt}
\begin{center}
\begin{tabular}{@{}lrrrrrrrrr@{}}
\toprule
& \multicolumn{4}{c}{brms} & \multicolumn{4}{c}{AutoBayes} & \\
\cmidrule(lr){2-5}\cmidrule(lr){6-9}& warm-up & sampling & depth & step & warm-up & sampling & depth & step & ratio \\
\midrule
B1 & 98,770 & 61,840 & 4.01 & 0.247 & 95,315 & 60,000 & 4.00 & 0.260 & 1.03 \\
B2 & 221,936 & 124,000 & 5.00 & 0.135 & 217,750 & 124,000 & 5.00 & 0.129 & 1.01 \\
B3 & 87,144 & 61,720 & 3.98 & 0.240 & 84,526 & 61,272 & 4.02 & 0.227 & 1.02 \\
B4 & 113,190 & 90,016 & 4.29 & 0.179 & 111,279 & 68,112 & 4.12 & 0.206 & 1.13 \\
B5 & 98,806 & 68,624 & 4.06 & 0.186 & 102,278 & 73,592 & 4.20 & 0.176 & 0.95 \\
B6 & 92,569 & 60,064 & 3.99 & 0.269 & 89,590 & 59,984 & 4.00 & 0.263 & 1.02 \\
B7 & 73,807 & 59,216 & 3.87 & 0.300 & 70,200 & 58,360 & 3.94 & 0.301 & 1.03 \\
B8 & 530,028 & 499,873 & 6.94 & 0.034 & 537,367 & 492,265 & 6.94 & 0.035 & 1.00 \\
B9 & 101,875 & 78,208 & 4.14 & 0.167 & 100,612 & 77,552 & 4.26 & 0.165 & 1.01 \\
B10 & 105,127 & 78,304 & 4.11 & 0.189 & 99,732 & 77,864 & 4.27 & 0.196 & 1.03 \\
B11 & 34,102 & 29,820 & 2.93 & 0.495 & 33,580 & 28,948 & 3.01 & 0.486 & 1.02 \\
B12 & 1,016,378 & 1,005,531 & 7.96 & 0.019 & 989,870 & 1,019,054 & 8.00 & 0.018 & 1.01 \\
B13 & 23,095 & 22,850 & 2.58 & 0.663 & 24,532 & 22,270 & 2.64 & 0.685 & 0.98 \\
B14 & 172,782 & 109,968 & 4.34 & 0.119 & 182,541 & 113,000 & 4.69 & 0.111 & 0.96 \\
B15 & 127,644 & 85,680 & 4.20 & 0.205 & 122,077 & 81,024 & 4.33 & 0.202 & 1.05 \\
B16 & 143,761 & 80,632 & 4.14 & 0.177 & 136,401 & 88,744 & 4.43 & 0.169 & 1.00 \\
B17 & 94,509 & 66,336 & 4.02 & 0.195 & 87,164 & 61,760 & 4.03 & 0.227 & 1.08 \\
B18 & 114,805 & 100,160 & 4.39 & 0.145 & 110,730 & 92,320 & 4.47 & 0.150 & 1.06 \\
\bottomrule
\end{tabular}\end{center}
\caption{\small Sampler work: leapfrog steps summed over the chains, for warm-up and for sampling; the mean tree depth while sampling and the final step size, averaged over the chains. The last column is brms's steps over AutoBayes's.}\label{tab:work}
\end{table}

\begin{table}[H]
\footnotesize
\setlength{\tabcolsep}{4pt}
\begin{center}
\begin{tabular}{@{}lrrrrrrrrrr@{}}
\toprule
& & \multicolumn{4}{c}{brms} & \multicolumn{4}{c}{AutoBayes} & \\
\cmidrule(lr){3-6}\cmidrule(lr){7-10}& coords & host & direct & null & used & host & null & loop & compile & ratio \\
\midrule
B1 & 105 & 109.0 & 103.0 & 5.0 & 98.0 & 109.5 & 7.1 & 52.7 & 0.16 & 1.86 \\
B2 & 407 & 809.2 & 805.1 & 5.0 & 800.1 & 250.4 & 7.7 & 110.6 & 0.16 & 7.24 \\
B3 & 36 & 770.1 & 761.0 & 5.0 & 756.0 & 167.8 & 7.3 & 102.1 & 0.15 & 7.40 \\
B4 & 44 & 253.9 & 242.0 & 5.0 & 237.0 & 122.4 & 7.2 & 73.6 & 0.13 & 3.22 \\
B5 & 44 & 129.0 & 123.0 & 5.0 & 118.0 & 140.1 & 6.9 & 77.9 & 0.16 & 1.51 \\
B6 & 45 & 381.0 & 376.0 & 5.0 & 371.0 & 282.8 & 7.0 & 290.4 & 0.28 & 1.28 \\
B7 & 36 & 473.0 & 467.1 & 5.0 & 462.1 & 213.2 & 6.9 & 145.2 & 0.13 & 3.18 \\
B8 & 44 & 16.0 & 11.0 & 5.0 & 6.0 & 87.0 & 7.1 & 13.9 & 0.19 & 0.43 \\
B9 & 44 & 474.6 & 476.8 & 5.0 & 471.8 & 362.1 & 7.0 & 285.0 & 0.28 & 1.66 \\
B10 & 34 & 55.1 & 50.1 & 5.0 & 45.1 & 52.0 & 7.3 & 4.6 & 0.10 & 9.73 \\
B11 & 7 & 221.0 & 212.0 & 5.0 & 206.9 & 92.2 & 7.2 & 32.8 & 0.20 & 6.31 \\
B12 & 204 & 922.0 & 908.1 & 5.0 & 903.1 & 1,287 & 7.0 & 952.0 & 0.20 & 0.95 \\
B13 & 4 & 158.1 & 150.9 & 5.0 & 145.9 & 63.3 & 6.8 & 12.1 & 0.10 & 12.09 \\
B14 & 24 & 78.0 & 73.9 & 5.0 & 68.9 & 53.0 & 7.2 & 5.9 & 0.09 & 11.65 \\
B15 & 69 & 227.0 & 219.1 & 5.0 & 214.1 & 81.8 & 8.4 & 12.1 & 0.19 & 17.74 \\
B16 & 38 & 109.0 & 99.9 & 5.0 & 94.9 & 72.7 & 7.7 & 11.6 & 0.16 & 8.20 \\
B17 & 44 & 139.0 & 134.0 & 5.0 & 129.0 & 180.4 & 7.0 & 117.0 & 0.18 & 1.10 \\
B18 & 45 & 145.0 & 143.1 & 5.0 & 138.0 & 80.8 & 7.0 & 22.5 & 0.10 & 6.14 \\
\bottomrule
\end{tabular}\end{center}
\caption{\small The cost of one log density and gradient, microseconds, at a posterior draw on the unconstrained scale (the number of coordinates first). brms: \code{rstan::grad\_log\_prob} (\emph{host}), the model instance's method called directly (\emph{direct}), a call that does no work (\emph{null}), and \emph{direct} less \emph{null}, the estimate used below. AutoBayes: the jitted potential and gradient called from Python (\emph{host}), a jitted call that does no work (\emph{null}), and the cost per evaluation inside one compiled loop (\emph{loop}), the estimate used below; \emph{compile} is seconds to trace, lower and compile it. Medians of 2,000 calls after a first one.}\label{tab:grad}
\end{table}

\section{Where the sampling time goes}\label{sec:decomp}

For each side, $T$ is the samplers' execution time (warm-up and sampling of every chain, compilation excluded), $S$ the leapfrog steps summed over the chains, $g$ the cost of one log density and gradient inside the sampler (Table~\ref{tab:grad}) and $C$ the CPU seconds spent over $T$. With $O = C/(Sg)$, the CPU cost of one step counted in gradients (1 if a step costs exactly a gradient), and $U = C/T$, the cores kept busy, $T = S g O / U$ holds exactly, so brms's $T$ over AutoBayes's is the product
\[ \frac{T_b}{T_a} = \underbrace{\frac{S_b}{S_a}}_{\text{steps}} \cdot \underbrace{\frac{g_b}{g_a}}_{\text{gradient}} \cdot \underbrace{\frac{O_b}{O_a}}_{\text{overhead}} \cdot \underbrace{\frac{U_a}{U_b}}_{\text{parallelism}}. \]
\textbf{A worked example, B11} (the median sampler speed-up). brms took $T = 3.97$ s and 63,922 leapfrog steps over all chains; at $g = 206.9\,\mu$s a gradient those steps alone are $S g = 13.23$ CPU seconds. Its chains used $C = 15.93$ CPU seconds, so each step cost $O = C/(Sg) = 1.20$ gradients and $U = C/T = 4.01$ cores were busy on average. AutoBayes took $T = 0.98$ s and 62,528 leapfrog steps over all chains; at $g = 32.8\,\mu$s a gradient those steps alone are $S g = 2.05$ CPU seconds. Its chains used $C = 4.48$ CPU seconds, so each step cost $O = C/(Sg) = 2.18$ gradients and $U = C/T = 4.56$ cores were busy on average. Then $3.97/0.98 = 4.05$, which is 1.02 (steps) times 6.31 (gradient) times 0.55 (overhead) times 1.14 (parallelism).

\begin{table}[H]
\footnotesize
\setlength{\tabcolsep}{4pt}
\begin{center}
\begin{tabular}{@{}lrrrrrrrrrrr@{}}
\toprule
& \multicolumn{2}{c}{$T$ (s)} & & \multicolumn{4}{c}{factor} & \multicolumn{2}{c}{$O$} & \multicolumn{2}{c}{$U$} \\
\cmidrule(lr){2-3}\cmidrule(lr){5-8}\cmidrule(lr){9-10}\cmidrule(lr){11-12}& brms & AutoBayes & ratio & steps & gradient & overhead & parallel & brms & AutoBayes & brms & AutoBayes \\
\midrule
B1 & 4.5 & 2.5 & 1.78 & 1.03 & 1.86 & 0.92 & 1.01 & 1.12 & 1.22 & 3.90 & 3.94 \\
B2 & 74.9 & 17.4 & 4.31 & 1.01 & 7.24 & 0.35 & 1.66 & 1.03 & 2.92 & 3.82 & 6.34 \\
B3 & 32.4 & 4.9 & 6.64 & 1.02 & 7.40 & 0.85 & 1.04 & 1.08 & 1.27 & 3.75 & 3.88 \\
B4 & 14.6 & 4.1 & 3.61 & 1.13 & 3.22 & 1.01 & 0.98 & 1.11 & 1.09 & 3.65 & 3.56 \\
B5 & 5.9 & 4.3 & 1.37 & 0.95 & 1.51 & 0.99 & 0.96 & 1.11 & 1.12 & 3.74 & 3.57 \\
B6 & 15.9 & 12.7 & 1.25 & 1.02 & 1.28 & 0.97 & 0.99 & 1.08 & 1.12 & 3.86 & 3.83 \\
B7 & 17.1 & 5.6 & 3.08 & 1.03 & 3.18 & 0.94 & 1.00 & 1.09 & 1.16 & 3.91 & 3.91 \\
B8 & 1.7 & 3.0 & 0.56 & 1.00 & 0.43 & 1.44 & 0.91 & 1.17 & 0.81 & 4.25 & 3.85 \\
B9 & 23.7 & 20.4 & 1.16 & 1.01 & 1.66 & 0.41 & 1.69 & 1.07 & 2.61 & 3.84 & 6.48 \\
B10 & 2.3 & 0.47 & 4.86 & 1.03 & 9.73 & 0.54 & 0.89 & 1.12 & 2.08 & 4.10 & 3.66 \\
B11 & 4.0 & 0.98 & 4.05 & 1.02 & 6.31 & 0.55 & 1.14 & 1.20 & 2.18 & 4.01 & 4.56 \\
B12 & 468.8 & 617.9 & 0.76 & 1.01 & 0.95 & 0.49 & 1.61 & 1.00 & 2.02 & 3.89 & 6.25 \\
B13 & 2.1 & 0.22 & 9.79 & 0.98 & 12.09 & 0.89 & 0.93 & 1.32 & 1.48 & 4.17 & 3.88 \\
B14 & 5.7 & 0.83 & 6.89 & 0.96 & 11.65 & 0.65 & 0.95 & 1.12 & 1.72 & 3.82 & 3.64 \\
B15 & 13.1 & 1.00 & 13.16 & 1.05 & 17.74 & 0.73 & 0.96 & 1.09 & 1.49 & 3.79 & 3.65 \\
B16 & 6.5 & 0.98 & 6.59 & 1.00 & 8.20 & 0.79 & 1.02 & 1.10 & 1.40 & 3.62 & 3.71 \\
B17 & 5.6 & 4.9 & 1.13 & 1.08 & 1.10 & 1.02 & 0.93 & 1.11 & 1.08 & 4.13 & 3.84 \\
B18 & 8.5 & 1.6 & 5.43 & 1.06 & 6.14 & 0.88 & 0.95 & 1.07 & 1.22 & 3.73 & 3.53 \\
\bottomrule
\end{tabular}\end{center}
\caption{\small The four factors of brms's sampler time over AutoBayes's, and each side's $O$ and $U$. A factor above 1 favours AutoBayes.}\label{tab:decomp}
\end{table}

\noindent Model by model, the largest factor:
\begin{itemize}
\item \textbf{B1} (Bernoulli, \code{(1 + x | g)}): sampler 1.78 times faster in AutoBayes; largest factor gradient (1.86, cheaper gradients).
\item \textbf{B2} (Gaussian, \code{(1 + x | g)}): sampler 4.31 times faster in AutoBayes; largest factor gradient (7.24, cheaper gradients).
\item \textbf{B3} (cumulative, five categories): sampler 6.64 times faster in AutoBayes; largest factor gradient (7.40, cheaper gradients).
\item \textbf{B4} (zero-inflated Poisson): sampler 3.61 times faster in AutoBayes; largest factor gradient (3.22, cheaper gradients).
\item \textbf{B5} (Weibull): sampler 1.37 times faster in AutoBayes; largest factor gradient (1.51, cheaper gradients).
\item \textbf{B6} (zero-inflated neg.\ binomial): sampler 1.25 times faster in AutoBayes; largest factor gradient (1.28, cheaper gradients).
\item \textbf{B7} (sratio, five categories): sampler 3.08 times faster in AutoBayes; largest factor gradient (3.18, cheaper gradients).
\item \textbf{B8} (Gaussian, horseshoe on 20): sampler 0.56 times faster in AutoBayes; largest factor gradient (0.43, against AutoBayes).
\item \textbf{B9} (beta-binomial, trials): sampler 1.16 times faster in AutoBayes; largest factor overhead (0.41, against AutoBayes).
\item \textbf{B10} (Gaussian, \code{(1 | mm(g1, g2))}): sampler 4.86 times faster in AutoBayes; largest factor gradient (9.73, cheaper gradients).
\item \textbf{B11} (mixture of two Gaussians): sampler 4.05 times faster in AutoBayes; largest factor gradient (6.31, cheaper gradients).
\item \textbf{B12} (exact GP, \code{gp(x)}): sampler 0.76 times faster in AutoBayes; largest factor overhead (0.49, against AutoBayes).
\item \textbf{B13} (AR(1), \code{ar(t, g)}): sampler 9.79 times faster in AutoBayes; largest factor gradient (12.09, cheaper gradients).
\item \textbf{B14} (non-linear, \code{a * exp(b * x)}): sampler 6.89 times faster in AutoBayes; largest factor gradient (11.65, cheaper gradients).
\item \textbf{B15} (\code{mvbind(y1, y2)}, rescor): sampler 13.16 times faster in AutoBayes; largest factor gradient (17.74, cheaper gradients).
\item \textbf{B16} (monotonic \code{mo(edu)}): sampler 6.59 times faster in AutoBayes; largest factor gradient (8.20, cheaper gradients).
\item \textbf{B17} (Weibull, \code{cens()}): sampler 1.13 times faster in AutoBayes; largest factor gradient (1.10, cheaper gradients).
\item \textbf{B18} (distributional, \code{sigma \textasciitilde{} z}): sampler 5.43 times faster in AutoBayes; largest factor gradient (6.14, cheaper gradients).
\end{itemize}

\section{Is the faster sampler sampling as well?}

\begin{table}[H]
\small
\setlength{\tabcolsep}{4pt}
\begin{center}
\begin{tabular}{@{}lrrrrrrrrr@{}}
\toprule
& \multicolumn{2}{c}{min bulk ESS} & \multicolumn{2}{c}{min tail ESS} & \multicolumn{2}{c}{bulk ESS/s} & \multicolumn{2}{c}{divergent} & max $|z|$ \\
\cmidrule(lr){2-3}\cmidrule(lr){4-5}\cmidrule(lr){6-7}\cmidrule(lr){8-9}& AutoBayes & brms & AutoBayes & brms & AutoBayes & brms & AutoBayes & brms & \\
\midrule
B1 & 1,264 & 1,017 & 1,843 & 1,638 & 500 & 226 & 0 & 0 & 1.88 \\
B2 & 125 & 144 & 298 & 459 & 7 & 2 & 0 & 0 & 1.22 \\
B3 & 530 & 558 & 1,062 & 1,076 & 109 & 17 & 0 & 0 & 0.41 \\
B4 & 386 & 414 & 726 & 951 & 95 & 28 & 0 & 0 & 1.46 \\
B5 & 201 & 287 & 296 & 541 & 47 & 49 & 0 & 0 & 1.02 \\
B6 & 1,079 & 1,245 & 1,874 & 2,019 & 85 & 78 & 0 & 0 & 0.70 \\
B7 & 1,087 & 1,195 & 1,418 & 1,717 & 196 & 70 & 0 & 0 & 1.86 \\
B8 & 1,916 & 2,257 & 1,846 & 1,313 & 637 & 1,342 & 9 & 6 & 2.71 \\
B9 & 185 & 264 & 465 & 499 & 9 & 11 & 0 & 0 & 1.43 \\
B10 & 533 & 499 & 929 & 835 & 1,143 & 220 & 0 & 0 & 0.82 \\
B11 & 3,325 & 3,131 & 2,974 & 2,922 & 3,390 & 788 & 0 & 0 & 1.49 \\
B12 & 1,214 & 1,410 & 1,709 & 1,799 & 2 & 3 & 2 & 4 & 0.57 \\
B13 & 4,174 & 4,536 & 2,801 & 2,842 & 19,306 & 2,143 & 0 & 0 & 0.78 \\
B14 & 290 & 269 & 446 & 379 & 351 & 47 & 0 & 0 & 0.51 \\
B15 & 806 & 1,022 & 1,499 & 1,754 & 808 & 78 & 0 & 0 & 1.11 \\
B16 & 440 & 400 & 675 & 893 & 448 & 62 & 0 & 0 & 1.36 \\
B17 & 355 & 309 & 625 & 660 & 72 & 56 & 0 & 0 & 0.69 \\
B18 & 192 & 200 & 377 & 277 & 122 & 23 & 0 & 0 & 1.30 \\
\bottomrule
\end{tabular}\end{center}
\caption{\small Sampling efficiency and agreement: the smallest bulk and tail ESS over the parameters both libraries report (the same ArviZ estimators for both), the smallest bulk ESS per second of the samplers' execution $T$, the divergent transitions after warm-up, and the largest $|z|$, the difference of the posterior means over their combined Monte Carlo standard error.}\label{tab:eff}
\end{table}

\section{Post-processing}

\begin{table}[H]
\footnotesize
\setlength{\tabcolsep}{3pt}
\begin{center}
\begin{tabular}{@{}lccccccc@{}}
\toprule
& \code{posterior\_epred} & \code{log\_lik} & \code{loo} & \code{waic} & \code{bayes\_R2} & \code{hypothesis} & \code{cond.\_effects} \\
\midrule
B1 & 0.12 / 0.35 & 0.53 / 0.71 & 2.16 / 3.23 & 0.64 / 1.79 & 0.16 / 0.63 & 0.00 / 0.00 & 0.02 / 0.15 \\
B2 & 0.68 / 1.20 & 2.49 / 5.27 & 18.97 / 32.10 & 3.84 / 17.60 & 0.89 / 3.80 & 0.00 / 0.00 & 0.02 / 0.17 \\
B3 & 1.77 / 1.65 & 0.58 / 1.23 & 1.77 / 3.67 & 0.66 / 2.46 & 1.94 / 2.35 & \textit{ref.} / 0.00 & \textit{ref.} / 0.27 \\
B4 & 0.13 / 0.35 & 0.35 / 1.39 & 2.00 / 4.09 & 0.47 / 2.92 & 0.17 / 0.63 & 0.00 / 0.00 & 0.02 / 0.14 \\
B5 & 0.12 / 0.38 & 0.19 / 1.06 & 1.87 / 3.50 & 0.29 / 2.11 & 0.16 / 0.68 & 0.00 / 0.00 & 0.02 / 0.15 \\
B6 & 0.14 / 0.36 & 1.00 / 1.80 & 2.65 / 4.26 & 1.11 / 3.14 & 0.18 / 0.66 & 0.00 / 0.00 & 0.02 / 0.15 \\
B7 & 1.31 / 9.78 & 1.22 / 1.95 & 2.35 / 4.40 & 1.28 / 3.03 & 1.44 / 10.09 & \textit{ref.} / 0.00 & \textit{ref.} / 0.76 \\
B8 & 0.04 / 0.06 & 0.03 / 0.10 & 0.33 / 0.58 & 0.05 / 0.35 & 0.05 / 0.13 & 0.00 / 0.00 & 0.03 / 0.20 \\
B9 & 0.16 / 0.41 & 1.48 / \textit{ref.} & 3.13 / \textit{ref.} & 1.59 / \textit{ref.} & 0.19 / 0.69 & 0.00 / 0.00 & \textit{ref.} / 0.15 \\
B10 & 0.05 / 0.10 & 0.08 / 0.21 & 0.82 / 1.61 & 0.14 / 0.99 & 0.07 / 0.21 & 0.00 / 0.00 & 0.02 / 0.15 \\
B11 & 0.05 / 0.12 & 0.29 / 0.60 & 1.00 / 2.18 & 0.32 / 1.40 & 0.06 / 0.33 & \textit{ref.} / 0.00 & \textit{ref.} / 0.15 \\
B12 & 2.66 / 3.58 & 2.00 / 3.71 & 2.53 / 3.97 & 2.55 / 3.79 & 1.77 / 3.57 & 0.00 / 0.00 & \textit{ref.} / 13.32 \\
B13 & 0.17 / 0.79 & 0.23 / 1.06 & 1.90 / 3.74 & 0.33 / 2.23 & 0.20 / 1.13 & 0.00 / 0.00 & \textit{ref.} / 0.15 \\
B14 & 0.06 / 0.32 & 0.09 / 0.43 & 0.86 / 1.62 & 0.15 / 1.20 & 0.08 / 0.44 & \textit{ref.} / 0.00 & \textit{ref.} / 0.16 \\
B15 & 0.10 / 0.55 & \textit{ref.} / 62.95 & \textit{ref.} / 64.63 & \textit{ref.} / 64.96 & \textit{ref.} / 0.71 & \textit{ref.} / 0.00 & \textit{ref.} / 0.47 \\
B16 & 0.09 / 0.34 & 0.13 / 0.68 & 1.30 / 2.43 & 0.21 / 1.61 & 0.11 / 0.66 & 0.00 / 0.00 & \textit{ref.} / 0.18 \\
B17 & 0.13 / 0.33 & 0.51 / 0.89 & 2.21 / 3.41 & 0.58 / 2.26 & 0.17 / 0.67 & 0.00 / 0.00 & \textit{ref.} / 0.15 \\
B18 & 0.09 / 0.41 & 0.23 / 0.69 & 1.96 / 3.32 & 0.33 / 2.17 & 0.13 / 0.80 & 0.00 / 0.00 & \textit{ref.} / 0.18 \\
\bottomrule
\end{tabular}\end{center}
\caption{\small Post-processing, seconds, AutoBayes / brms: the median of three calls after a first one. \textit{ref.}: the function refused the model (the message is in the JSON).}\label{tab:post}
\end{table}

\section{Against the previous run}

\begin{table}[H]
\small
\setlength{\tabcolsep}{4pt}
\begin{center}
\begin{tabular}{@{}lrrrr@{}}
\toprule
& \multicolumn{2}{c}{brms} & \multicolumn{2}{c}{AutoBayes} \\
\cmidrule(lr){2-3}\cmidrule(lr){4-5}& before & now & before & now \\
\midrule
B1 & 42.0 & 36.8 & 5.3 & 7.0 \\
B2 & 122.5 & 105.9 & 29.6 & 30.5 \\
B3 & 60.4 & 58.9 & 7.1 & 6.7 \\
B4 & 41.2 & 40.6 & 5.9 & 6.0 \\
B5 & 33.5 & 32.7 & 6.0 & 6.0 \\
B6 & 43.0 & 42.7 & 14.7 & 15.2 \\
B7 & 43.6 & 43.5 & 7.1 & 7.1 \\
B8 & 29.6 & 27.9 & 5.3 & 5.1 \\
B9 & 52.1 & 49.1 & 22.3 & 22.6 \\
B10 & 30.0 & 27.6 & 2.4 & 1.9 \\
\bottomrule
\end{tabular}\end{center}
\caption{\small End-to-end seconds in the previous run and in this one, for the models both have.}\label{tab:prev}
\end{table}

\section{What these numbers do not show}

\begin{itemize}
\item \textbf{One run each}, on one machine whose load is not controlled: a single number may be off by 10--30\%.
\item \textbf{Step counts differ for reasons beyond the gradient.} Stan and NumPyro both run NUTS with a diagonal metric
  and dual-averaging step-size adaptation, but their warm-up schedules, windows and step-size searches are not the same,
  and they start from different points, so they end on different step sizes and tree depths. The steps factor measures
  that difference; it says nothing about either library's gradient. Neither count includes the gradients spent searching
  for an initial step size.
\item \textbf{The gradient costs are estimates.} The host-call medians include a call overhead the samplers do not pay
  (for a jitted JAX function on the CPU, tens of microseconds on a busy machine); the factors use AutoBayes's cost inside
  a compiled loop and rstan's direct call less a call that does no work. brms's point is a draw of a short extra chain
  (brms renames its fit's parameters, so the fit's draws cannot be mapped back to Stan's coordinates); AutoBayes's is the
  last draw of its fit. For these models the cost hardly depends on the point.
\item \textbf{CPU time and parallelism.} AutoBayes's $C$ is the whole process's CPU time over the run's execution
  (XLA's threads included); brms's is the CPU time of the forked chain processes over \code{update()}, which includes
  each child's small R-side setup. $U$ above the number of chains means a chain used more than one core.
\item \textbf{Compilation is measured from the timeline.} AutoBayes's sampler compilation is the time from the start of
  the NumPyro run to the end of its largest XLA compilation; the part JAX reports as tracing, lowering and compiling is
  shown beside it. brms's codegen is a warm re-run of \code{stancode()} and \code{standata()}.
\item \textbf{Repeated fits.} brms compiles once per model and can refit new data without recompiling; for many fits of
  one model only the samplers' time matters.
\item \textbf{Other setups.} rstan only (cmdstanr is not installed), no within-chain threading, the CPU only.
\end{itemize}

\section{Reproducing}

\begin{verbatim}
python examples/brms_parity/benchmark_data.py DATA
Rscript examples/brms_parity/benchmark_brms.R DATA OUT
python examples/brms_parity/benchmark_autobayes.py DATA OUT
python examples/brms_parity/compare.py OUT
python examples/brms_parity/report.py OUT/comparison-DATE.json --pdf
\end{verbatim}

\end{document}
